% ============================================================================
% Patch for Section 2 (and Summary) so that the paper states what was actually
% fitted, namely delta Phi.  K. Yasuda, 2026-08-20.
% Replaces the paragraph currently beginning "The rest of this section needs to
% be changed based on delta Phi..." (the TODO note is deleted).
% ============================================================================

% ---------------------------------------------------------------- Section 2 --
% [REPLACE the paragraph "To work with a function whose value is defined at most
%  energies, we fit the data with the integrated flux Phi(E_nu). Since any
%  likelihood ... This is what we find in the present application."  by:]

To work with a function whose value is defined at most energies, we fit the data with an
integrated flux rather than with the differential one. The reactor antineutrino spectrum
above $2$~MeV is known from measurements above the IBD threshold and we take it as a fixed
input, so only the flux below $2$~MeV is to be determined. As we show in Sec.~4, once the
contribution to the rate of the neutrinos with $E_\nu > 2$~MeV is separated, the measured
rate depends linearly not on $\Phi(E_\nu)$ but on the flux integrated between $E_\nu$ and
$2$~MeV,
\begin{equation}
\delta\Phi(E_\nu)\;\equiv\;\Phi(E_\nu)-\Phi(2\,\mathrm{MeV})
\;=\;\int_{E_\nu}^{2\,\mathrm{MeV}}\! d\tilde E_\nu\,\phi(\tilde E_\nu)\,,
\end{equation}
defined for $E_\nu^{\min}\le E_\nu\le 2$~MeV, where $E_\nu^{\min}$ is the smallest neutrino
energy that can produce a recoil above the detector threshold (see Sec.~5). It is
$\delta\Phi(E_\nu)$ that we determine by fitting the data. The integrated flux is then
obtained as $\Phi(E_\nu)=\delta\Phi(E_\nu)+\Phi(2\,\mathrm{MeV})$, which is the quantity we
display in Figs.~7--11.

Because $\delta\Phi$ and $\Phi$ differ by a constant, $\delta\Phi$ is non-negative and
non-increasing, and it is piecewise constant with a downward step wherever $\Phi$ has one.
The results just quoted therefore hold for $\delta\Phi$ without modification: any likelihood
can be maximized with a $\delta\Phi$ of the form of Eq.~(2.5) (with the sum restricted to
steps below $2$~MeV), so that the problem of finding the function $\delta\Phi(E_\nu)$ becomes
one of finding at most $2(d-1)$ parameters. Although any likelihood can be maximized with
a function of this form, there may be other functions that maximize it as well, in which case
there is a ``degeneracy band'' around the best fit, as explained in Ref.~[8]. This is what we
find in the present application.

% [In the SAME section, the paragraph describing the numerical ansatz currently
%  reads "... in each of which Phi(E_nu) takes a constant value, Phi_j, as shown
%  in Fig. 4. With this ansatz for Phi(E_nu), chi^2 becomes a function of the
%  N_int parameters Phi_j. The best-fit Phi_j values are determined by minimizing
%  chi^2 while imposing the physical requirement that Phi(E_nu) be non-increasing,
%  i.e., that the Phi_j values cannot increase with increasing j. In the best-fit
%  function, many adjacent intervals have identical Phi_j values, ..."
%  REPLACE Phi -> delta Phi throughout it:]

Applied to neutrino scattering, the method proceeds as follows. Divide the neutrino energy
range of interest into a large number $N_{\rm int}$ of equally spaced intervals
$j=1\dots N_{\rm int}$, in each of which $\delta\Phi(E_\nu)$ takes a constant value,
$\delta\Phi_j$, as shown in Fig.~4. With this ansatz for $\delta\Phi(E_\nu)$,
$\chi^2_{\rm Neyman}$ becomes a function of the $N_{\rm int}$ parameters $\delta\Phi_j$. The
best-fit $\delta\Phi_j$ values are determined by minimizing $\chi^2_{\rm Neyman}$ while
imposing the physical requirement that $\delta\Phi(E_\nu)$ be non-increasing, i.e., that the
$\delta\Phi_j$ values cannot increase with increasing $j$. In the best-fit function, many
adjacent intervals have identical $\delta\Phi_j$ values, forming a constant segment, resulting
in no more than $d-1$ downward steps.

% ------------------------------------------------------------------ Summary --
% [In Sec. 7, two sentences still describe the fit as being of the integrated
%  flux.  Suggested replacements:]

% "The resulting integrated neutrino spectrum as a function of the lowest energy in each bin
%  is thus piecewise constant ... The location and height of each step of the integrated flux
%  need to be determined by likelihood maximization"
%   ->
The function we unfold is $\delta\Phi(E_\nu)$, the neutrino flux integrated between $E_\nu$
and $2$~MeV, which is piecewise constant with at most $d-1$ downward steps; the integrated
flux above $E_\nu$ follows by adding the known constant $\Phi(2\,\mathrm{MeV})$. The location
and height of each step need to be determined by likelihood maximization, which is a total of
at most $2(d-1)$ parameters.

% "... dividing the neutrino energy range of interest in a very large number N_int of equally
%  spaced intervals, in each of which the integrated flux takes a unique value, ..."
%   ->
% "... in each of which $\delta\Phi$ takes a unique value, ..."

% "A limitation of our current approach is that it fits and defines confidence bands for the
%  integrated flux."
%   ->
% "A limitation of our current approach is that it fits and defines confidence bands for the
%  integrated flux, and not for the differential one."   (unchanged in meaning; optional)

% ----------------------------------------------------------------- Abstract --
% Optional, for accuracy: "... a pointwise confidence band for the integrated neutrino flux"
% -> "... for the integrated neutrino flux below 2 MeV".
